\section{The passages and prayers}

The biblical text here is the public-domain Douay--Rheims Bible, Challoner revision (Gutenberg 1581). At Mass the approved United States \textit{Lectionary for Mass} supplies the proclaimed English. Psalms are identified by their Hebrew number and then their Vulgate number. The verse numbers beside the study text follow Challoner. The \textit{Roman Missal}, Third Edition for the United States, gives the antiphons and prayers.

\subsection{Entrance antiphon}
\label{proper-entrance}
\textit{Ad te levavi}, First Sunday of Advent, Entrance Antiphon (U.S. Missal, p.~139). The Missal adapts Psalm 25 (24):1--3. Here is the biblical source, beginning after the superscription. The adaptation itself is reserved.
\begin{quote}
\textsuperscript{1} To thee, O Lord, have I lifted up my soul.
\textsuperscript{2} In thee, O my God, I put my trust; let me not be ashamed.
\textsuperscript{3} Neither let my enemies laugh at me: for none of them that wait on thee shall be confounded.
\end{quote}

\subsection{Collect}
\label{proper-collect}
\textit{Da, quaesumus, omnipotens Deus}, First Sunday of Advent (U.S. Missal, p.~139). The petition concerns the will to meet Christ and the deeds of those who stand at his right hand. The approved English and modern Latin prayer are reserved.

\subsection{First reading}
\label{proper-first-reading}
Isaiah 63:16b--17, 19b and 64:2--7 in the U.S. Lectionary, vol.~I, no.~2. The chapter division in Challoner differs. Begin at the second address to the Father in 63:16 and resume after 63:17 at 64:1. Skip native 64:2 and resume at 64:3. Challoner's wonderful deeds in 64:3 look forward before recalling a descent and the melting mountains. The U.S. proclamation recalls past wonders and does not repeat that close. At 64:5 Challoner speaks of God having met a just person and of others remembering God. The U.S. Isaiah 64:4 wishes for God to meet the speakers themselves doing right and remembering him. Challoner also affirms future salvation at the end of that verse, where the U.S. passage proceeds from sin to uncleanness. In Challoner 64:7 the speakers describe themselves as crushed under their own sin. In the U.S. text God has given them over to their guilt. The potter in the U.S. close is called the maker here.
\begin{quote}
\textsuperscript{16} thou, O Lord, art our father, our redeemer, from everlasting is thy name.
\textsuperscript{17} Why hast thou made us to err, O Lord, from thy ways: why hast thou hardened our heart, that we should not fear thee? return for the sake of thy servants, the tribes of thy inheritance.

\textsuperscript{1} O that thou wouldst rend the heavens, and wouldst come down: the mountains would melt away at thy presence.

\textsuperscript{3} When thou shalt do wonderful things, we shall not bear them: thou didst come down, and at thy presence the mountains melted away.
\textsuperscript{4} From the beginning of the world they have not heard, nor perceived with the ears: the eye hath not seen, O God, besides thee, what things thou hast prepared for them that wait for thee.
\textsuperscript{5} Thou hast met him that rejoiceth, and doth justice: in thy ways they shall remember thee: behold thou art angry, and we have sinned: in them we have been always, and we shall be saved.
\textsuperscript{6} And we are all become as one unclean, and all our justices as the rag of a menstruous woman: and we have all fallen as a leaf, and our iniquities, like the wind, have taken us away.
\textsuperscript{7} There is none that calleth upon thy name: that riseth up, and taketh hold of thee: thou hast hid thy face from us, and hast crushed us in the hand of our iniquity.
\textsuperscript{8} And now, O Lord, thou art our father, and we are clay: and thou art our maker, and we all are the works of thy hands.
\end{quote}

\subsection{Responsorial psalm}
\label{proper-psalm}
Psalm 80 (79):2ac, 3b, 15--16, 18--19. The U.S. heading gives 2--3 for the first strophe. The clause naming Joseph and the names of the three tribes lie outside the selection. The response is based on v.~4. Here is that biblical verse followed by the three strophes in study English. The English response sung at Mass has its own opening address and is reserved.
\begin{quote}
\textsuperscript{4} Convert us, O God: and shew us thy face, and we shall be saved.

\textsuperscript{2} Give ear, O thou that rulest Israel: Thou that sittest upon the cherubims, shine forth
\textsuperscript{3} Stir up thy might, and come to save us.

\textsuperscript{15} Turn again, O God of hosts, look down from heaven, and see, and visit this vineyard:
\textsuperscript{16} And perfect the same which thy right hand hath planted: and upon the son of man whom thou hast confirmed for thyself.

\textsuperscript{18} Let thy hand be upon the man of thy right hand: and upon the son of man whom thou hast confirmed for thyself.
\textsuperscript{19} And we depart not from thee, thou shalt quicken us: and we will call upon thy name.
\end{quote}
The seasonal sung choices and their passages are given under Liturgical Resolution.

\subsection{Second reading}
\label{proper-second-reading}
1 Corinthians 1:3--9. The U.S. Lectionary introduces the reading with its own address to the hearers.
\begin{quote}
\textsuperscript{3} Grace to you and peace, from God our father and from the Lord Jesus Christ.
\textsuperscript{4} I give thanks to my God always for you, for the grace of God that is given you in Christ Jesus:
\textsuperscript{5} That in all things you are made rich in him, in all utterance and in all knowledge;
\textsuperscript{6} As the testimony of Christ was confirmed in you,
\textsuperscript{7} So that nothing is wanting to you in any grace, waiting for the manifestation of our Lord Jesus Christ.
\textsuperscript{8} Who also will confirm you unto the end without crime, in the day of the coming of our Lord Jesus Christ.
\textsuperscript{9} God is faithful: by whom you are called unto the fellowship of his Son, Jesus Christ our Lord.
\end{quote}

\subsection{Alleluia verse}
\label{proper-acclamation}
Psalm 85 (84):8, the biblical verse for the Alleluia (U.S. Lectionary, no.~2, and Latin \textit{Ordo lectionum Missae}, no.~2 referring to no.~1). The protected liturgical English is omitted.
\begin{quote}
\textsuperscript{8} Shew us, O Lord, thy mercy; and grant us thy salvation.
\end{quote}

\subsection{Gospel}
\label{proper-gospel}
Mark 13:33--37. The Lectionary identifies Jesus as the speaker for those who hear the passage on its own. Its opening command lacks the words \emph{and pray} in Challoner v.~33. The U.S. rendering of v.~34 explicitly gives each servant work of his own as well as assigning the porter to keep watch.
\begin{quote}
\textsuperscript{33} Take ye heed, watch and pray. For ye know not when the time is.
\textsuperscript{34} Even as a man who, going into a far country, left his house and gave authority to his servants over every work and commanded the porter to watch.
\textsuperscript{35} Watch ye therefore (for you know not when the lord of the house cometh, at even, or at midnight, or at the cock crowing, or in the morning):
\textsuperscript{36} Lest coming on a sudden, he find you sleeping.
\textsuperscript{37} And what I say to you, I say to all: Watch.
\end{quote}

\subsection{Prayer over the Offerings}
\label{proper-offerings}
\textit{Suscipe, quaesumus, Domine}, First Sunday of Advent (U.S. Missal, p.~139). The gifts on the altar came into the offerers' hands through an earlier gift from God. The approved English and modern Latin prayer are reserved.

\subsection{Preface I of Advent}
\label{proper-preface}
\textit{De duobus adventibus Christi}, on Christ's two comings (U.S. Missal, Order of Mass no.~33, p.~534). The priest gives thanks for Christ's coming in human lowliness and looks toward his appearing in glory. The approved English and modern Latin Preface are reserved.

\subsection{Communion antiphon}
\label{proper-communion}
\textit{Dominus dabit benignitatem}, First Sunday of Advent (U.S. Missal, p.~139). The biblical source is Psalm 85 (84):13. The Missal's own English is reserved.
\begin{quote}
\textsuperscript{13} For the Lord will give goodness: and our earth shall yield her fruit.
\end{quote}

\subsection{Prayer after Communion}
\label{proper-after-communion}
\textit{Prosint nobis}, First Sunday of Advent (U.S. Missal, p.~139). The petition concerns what those who have received Communion can learn as they handle things that pass away. The approved English and modern Latin prayer are reserved.

\subsection{Solemn blessing for Advent}
\label{proper-advent-blessing}
U.S. Missal, Solemn Blessings I.1, Advent, p.~674. The direction on p.~139 permits this option. Three invocations precede the final blessing. The approved English and modern Latin prayers are reserved.
